\subsection{到達段階と構造検証の範囲}
L0〜L3の累積達成と排他的な最終段階，baselineからの5×5遷移を固定CSVに保存した．L2とL3の件数は一致し，L2のみ達成は0件だったが，定義上の同一性を意味しない．統合のR・F・Dで，段階上昇したもののL3未達の部分改善は38・49・51件，最終段階の低下は27・25・31件である．段階は順序カテゴリであり，差を等間隔の点数として平均しない．

構造検証は量子ビット数と測定を除く回路命令数の範囲に限る．全10タスクで古典ビット数・測定数・測定の必須／禁止は未指定であり，ゲート種・接続・深さ・最適性や全入力での正しさを完全には検証しない．生成回路への構造要求と，評価用コピーへの測定追加を区別する．

\subsection{FとDの対応付き救出と層別結果}
FとDは各191/582件（32.8\%）を救出したが，共通救出149件，Fのみ42件，Dのみ42件，双方未達349件だった．同じ総数でも救出対象は異なる．表\ref{tab:discordant}は初期失敗類型別の対応を示す．importではDのみ18件，構文ではFのみ9件が観測された．保存済み履歴から，片方が成功する間にもう片方が同類型を繰り返す例や別類型へ移る例を確認した．分類の変化だけで具体的な修正原因は確定しない．
\begin{table}[t]
\centering
\caption{初期失敗類型別のF・D救出の対応}
\label{tab:discordant}
\setlength{\tabcolsep}{3pt}
\setlength{\tabcolsep}{3pt}
\begin{tabular}{l|rrrrr}\hline
類型 & 分母 & 共通 & Fのみ & Dのみ & 双方未達 \\\hline
import & 271 & 47 & 8 & 18 & 198 \\
回路構築 & 106 & 46 & 11 & 11 & 38 \\
インターフェース & 32 & 11 & 3 & 2 & 16 \\
構文 & 68 & 16 & 9 & 2 & 41 \\
出力不一致 & 102 & 29 & 11 & 9 & 53 \\
量子ビット数 & 3 & 0 & 0 & 0 & 3 \\
\hline
合計 & 582 & 149 & 42 & 42 & 349 \\\hline
\end{tabular}
\end{table}
モデル別ではG4がF 50/67・D 43/67，Q4がF 74/165・D 83/165，G3が双方0/198だった．易・中・難の救出はF 25/74・77/248・89/260，D 29/74・83/248・79/260．モデル・難度・seedブロック・出典別の内訳を保存し，全体の観測最多を一律の介入選択規則にしない．

\subsection{修正回数と追加生成コスト}
共有初期生成を除いたround 1以降の追加コストを集計した．未救出例の費用も含む追加tokens／救出はR 17,365.1，F 9,584.8，D 11,103.4，要求秒／救出は841.6，464.8，574.8だった．時間は保存されたLLM要求の経過時間で，評価処理の時間を含まない．未救出例が修正tokensの88.7\%，82.0\%，81.7\%を占めたが，部分改善もあるため無価値とは断定しない．

保存済み履歴を上限1・2・3回で区切った累積救出は，R 65→90→110件，F 119→166→191件，D 114→167→191件だった．上限変更の新規実験ではない．表\ref{tab:thirdround}は3回目の限界救出と当該roundだけの追加コストを示す．分母は2回目まで未達で実際に3回目を行った組数である．
\begin{table}[t]
\centering
\caption{修正3回目の救出と追加コスト}
\label{tab:thirdround}
\begin{tabular}{l|rrrr}\hline
条件 & 実施組 & 新規救出 & 追加tokens & 要求時間（h） \\\hline
R & 492 & 20 & 584,463 & 7.17 \\
F & 416 & 25 & 528,848 & 6.70 \\
D & 415 & 24 & 619,981 & 8.35 \\
\hline\end{tabular}
\end{table}
F 25件とD 24件の3回目救出は，共通3組を除く46初期組である．round 1・2が同じ失敗分類なら停止する候補は，Fで18/25件，Dで15/24件の救出を失う．初期から段階上昇した例だけ続ける候補も19/25件，21/24件を失う．これらは同じ履歴を用いた探索的集計であり，運用上の停止方策は未検証である．

\subsection{出力表現への感度と評価器の更新}
2026年10月14日にLLMを使わない確認用回路19例を実行した．正しいQPEでも未使用古典ビット3個により出力が1010から1010000となり，旧比較では不一致になった．測定先から量子ビット順へ戻す比較では，正しい10例中，旧比較で失敗した7例が成功した．誤った6例は不合格を維持し，曖昧な測定3例は拒否した．少数の制御例であり，あらゆる誤回路への保証ではない．

保存済み分布2,880 round行中619行で，タスクの量子ビット数と出力幅が異なった．共有初期生成の重複を含み，部分測定や回路自体の誤りもあり得るため，619行を独立生成数や表現だけによる失敗数とは扱わない．これらを一括して成功へ変更する根拠はない．

本番評価器へ\texttt{qubit\_readout\_v2}を採用し，未使用古典ビットを合算し，全量子ビットの一意な終端測定から分布を復元する仕様とした．部分測定・上書き・測定後の量子操作等は拒否し，新規manifestに版とhashを記録する．\textbf{本稿の5,000件は旧評価器の保存判定であり，新版による再評価結果ではない．} 実装時の257テスト成功は処理の検証であり，本稿の成功率が新版でも維持されることを示さない．
